\documentclass[11pt]{letter}
\usepackage[margin=1in]{geometry}
\usepackage{times}
\usepackage[hidelinks]{hyperref}

\signature{Zhouning}
\address{AlphaEarth-System Project \\ Inner Mongolia, China}

\begin{document}

\begin{letter}{%
  Editor-in-Chief \\
  \emph{Remote Sensing of Environment} \\
}

\opening{Dear Editor,}

We are pleased to submit our manuscript, ``How Well Do PEFT Methods Adapt Geospatial Foundation Models to Heterogeneous Inputs? A Systematic Evaluation on Prithvi-100M with Production Validation on a Chinese Domestic High-Resolution Corpus,'' for consideration as a research article in \emph{Remote Sensing of Environment}.

\textbf{Scope fit.} The manuscript reports a systematic empirical study of how parameter-efficient fine-tuning (PEFT) methods adapt the publicly available Prithvi-100M geospatial foundation model to heterogeneous remote sensing inputs, validates the laboratory benchmark on a 35{,}920-patch operational deployment over Linhe County, China, using 203 multi-sensor scenes (0.42--4.29\,m ground sampling distance, 22 satellite families, six calendar quarters in 2025) with Esri 10\,m global LULC as the supervisory signal, and finally replicates the same five-method protocol on the LoveDA urban/rural cross-domain benchmark to test an adapter-capacity hypothesis under explicit domain shift. This benchmark--deployment--cross-domain triple is, we believe, well aligned with RSE's emphasis on environmental applications grounded in rigorous remote sensing methodology rather than method papers in isolation.

\textbf{Principal contributions.}
\begin{enumerate}
    \item A controlled benchmark of five PEFT methods (linear probing, BitFit, LoRA at rank 8, Houlsby adapters, an input-stage GeoAdapter) across single-label classification (EuroSAT, 27 modality configurations), multi-label classification (BigEarthNet-S2), and semantic segmentation (LandCover.ai). 109 experiments in total.
    \item A diagnostic finding that standard module-wise LoRA can silently miss Prithvi-100M's PyTorch fused-attention projections because Q/K/V are packed into a single tensor. A split-QKV ablation isolates implementation failure from genuine method limitation.
    \item Production validation on the AlphaEarth-System platform: Houlsby adapters deliver $+0.116$ absolute mIoU ($+65.5\%$ relative) over linear probing on the Linhe LULC corpus under a geographic (scene-level) split protocol. The Houlsby-over-linear ordering is preserved across EuroSAT, BigEarthNet-S2, LandCover.ai, Linhe, and LoveDA within the tested Prithvi-100M setting.
    \item A complementary synthetic weak-label control showing that the PEFT delta on the same imagery is mediated by label source: brightness-thresholded weak labels yield only a $+0.017$ PEFT delta, while Esri-derived supervisory labels yield $+0.116$. This control is not an independent manual validation set, but it helps separate method behavior from label-source effects.
    \item A 30-run cross-domain replication on LoveDA (5 methods $\times$ 2 directions $\times$ 3 seeds) reporting evidence consistent with an adapter-capacity hypothesis near $10^6$ trainable parameters in this tested setting. Houlsby reaches $+0.117$ mIoU over linear probing under R$\rightarrow$U transfer; every PEFT method below the threshold (linear probing, BitFit, LoRA r=8/16, GeoAdapter) plateaus within $10^{-3}$ mIoU of the all-background baseline despite a 7--9$\times$ trainable-parameter advantage over linear probing. A per-class diagnostic at 5$\times$ learning rate and 60 epochs confirms that the smaller adapters recover at most a binary background-versus-foreground distinction rather than pure collapse or simple undertraining. We therefore present the capacity threshold as a Prithvi-100M/LoveDA diagnostic hypothesis, not as a general deployment rule.
\end{enumerate}

\textbf{Reproducibility.} All training code, configuration files, per-seed log files, and the AlphaEarth-System platform that orchestrates the Linhe pipeline are released under an open-source license. The Linhe patch index ($\sim$35\,k patches at $128\times128$, RGB) is shareable subject to source-imagery licensing; the supervisory labels are entirely derived from the publicly available Esri 10\,m LULC product.

\textbf{Confirmation.} This manuscript has not been published, accepted, or simultaneously submitted elsewhere. An earlier preprint version covering the public-benchmark portion only (Sections 1--8) was prepared for the CVPR 2026 EarthVision workshop but was not submitted following the workshop deadline; the present manuscript adds Section 9 (production validation on Linhe and cross-domain replication on LoveDA) and substantially revises the discussion. All experimental decisions and authorship are unchanged.

\textbf{Suggested reviewers.} We suggest, but defer to your judgment, reviewers with expertise in geospatial foundation models and remote sensing PEFT (e.g., authors of Prithvi-100M / SatMAE / Scale-MAE / SpectralGPT), Chinese-domestic high-resolution remote sensing benchmarks, and operational LULC mapping.

We appreciate your consideration of our manuscript.

\closing{Sincerely,}

\end{letter}
\end{document}
